\documentclass[11pt,a4paper]{article}
\usepackage[margin=1in]{geometry}
\usepackage{hyperref}
\usepackage{enumitem}
\usepackage{graphicx}
\usepackage{xcolor}

% -----------------------------------------------------
% bvraghav-tiet-ucs503p-article.sty
% -----------------------------------------------------

% Variable: Lab Instructor
\makeatletter \newcommand{\BvrSetLabInstructor}[1]{%
  \gdef\Bvr@LabInstructor{#1}}
% default
\newcommand{\Bvr@LabInstructor}{Dr. Raghav %
  B. Venkataramaiyer}
% public accessor
\newcommand{\BvrLabInstructorValue}{\Bvr@LabInstructor}
\makeatother

% Add subtitle
\usepackage{titling}
% -- Set title bold
\pretitle{\begin{center}\bfseries\fontsize{18bp}{18bp}\selectfont}
\makeatletter
\newcommand{\BvrSubtitle}[1]{%
  \posttitle{%
    \par\end{center}
    \begin{center}\large#1\end{center}
    \vskip 0.5em}%
}
\makeatother

% Hints in the section title(s)
\newcommand{\BvrHint}[1]{%
  {\normalfont\normalsize\itshape\hfill #1}}

% -----------------------------------------------------

\title{AI-Powered Student Workload Balancer}

\BvrSetLabInstructor{Dr. Nisha Thakur}

\BvrSubtitle{%
  Project Proposal for UCS503P  \\
  Submitted to: \BvrLabInstructorValue \\ \bfseries
  Thapar Institute of Engineering and Technology%
}

\author{
  Khushi, CSED%
  \thanks{Roll No: \texttt{1024030809}, %
    Email: \texttt{kkhushi1\_be24\@thapar.edu}}
  \and
  Umang Gautam, CSED%
  \thanks{Roll No: \texttt{1024030805}, %
    Email: \texttt{ugautam\_be24\@thapar.edu}}
  \and
  Vriti Mahajan, CSED%
  \thanks{Roll No: \texttt{1024030801}, %
    Email: \texttt{vmahajan1\_be24\@thapar.edu}}
  \and
  Shaurya Garg, CSED%
  \thanks{Roll No: \texttt{1024030800}, %
    Email: \texttt{sgarg6\_be24\@thapar.edu}}
}

\date{\today}

\begin{document}
\maketitle

\section{Higher-order goal%
  \BvrHint{\ldots secondary framing}}
This project aims to improve student time-management outcomes by
reducing the cognitive overhead of planning academic workload across
subjects, topics, and deadlines. While the broader goal is student
well-being and academic performance, the immediate engineering
objective is to translate \emph{adaptive workload planning} into a
measurable, deployable software solution.

\section{Time-to-value %
\BvrHint{\ldots primary heuristic}}
\begin{itemize}[leftmargin=0.7cm]
    \item We will deliver meaningful increments quickly by prototyping the core planning workflow (input subjects/hours/deadlines $\rightarrow$ generated plan) and validating it with users within a short iteration window.
    \item We will prioritize fast feedback over perfection by using weekly iterations (and targeting CI-backed automation early) to reduce release friction.
    \item Success is defined by what can be delivered and measured in the first iteration, then improved based on user feedback.
\end{itemize}

\section{Problem Statement}
Students juggling multiple subjects, assignments, and deadlines
struggle to allocate limited study time effectively. Common pain
points include:
\begin{itemize}[leftmargin=0.7cm]
    \item \textbf{Static planning:} manually made schedules quickly become outdated once a student falls behind or a deadline shifts.
    \item \textbf{No feedback loop:} most planning tools do not adapt based on actual progress or performance, leaving students to manually re-plan.
    \item \textbf{Uneven workload distribution:} without a systematic priority scheme, students over-invest in familiar subjects and under-invest in weak or high-stakes ones.
    \item \textbf{Opaque reasoning:} even when a tool suggests a plan, students rarely understand \emph{why} that plan was chosen, reducing trust and adoption.
\end{itemize}

\textbf{Why this matters now (contextual relevance):} With increasing academic load and limited study hours, even small improvements in workload distribution and re-planning speed can noticeably reduce student stress and missed deadlines.

\section{Proposed Solution %
\BvrHint{\ldots Software Proposition}}
\subsection{Overview}
Build a \textbf{web-based} ``AI-Powered Student Workload Balancer''
with the following components:
\begin{itemize}[leftmargin=0.7cm]
    \item \textbf{Student input portal:} students enter subjects, topics, available study hours, deadlines, and assignments.
    \item \textbf{Deterministic workload engine:} a rules-based backend service computes a personalized, prioritized study plan from the inputs and performance history.
    \item \textbf{Adaptive feedback loop:} the plan is continuously re-balanced based on actual progress and performance reported by the student.
    \item \textbf{Reasoning agent:} a LangGraph-based agent explains \emph{why} the plan looks the way it does, without altering the underlying business logic.
    \item \textbf{Progress dashboard:} students track plan adherence and upcoming priorities.
\end{itemize}

\subsection{Core workflow}
\begin{enumerate}[leftmargin=0.7cm]
    \item Student registers subjects, topics, available hours, and deadlines.
    \item The deterministic workload engine computes a prioritized study plan using a scoring algorithm (deadline proximity, topic weight, past performance).
    \item Student follows the plan and logs actual progress (completed/missed sessions, performance on topics).
    \item The engine re-balances the remaining plan based on this feedback; the agent explains the change in plain language on request.
\end{enumerate}

\subsection{Operational constraints for an educational setting}
\begin{itemize}[leftmargin=0.7cm]
    \item Keep the solution usable on low-to-mid range devices via responsive web design.
    \item Keep all scheduling logic deterministic and auditable; the AI agent may only reason over and explain outputs, never bypass business rules.
    \item Design for deployability using common web hosting and managed databases.
\end{itemize}

\section{Solution Approach %
\BvrHint{\ldots Engineering focus}}
This is an engineering course project, so we focus on a layered,
maintainable system architecture rather than novel algorithm
research.

\subsection{Deterministic engine + reasoning agent %
\BvrHint{\ldots non-research}}
Prioritization will be heuristic-based, not black-box:
\begin{itemize}[leftmargin=0.7cm]
    \item Compute a priority score per topic from deadline proximity, estimated effort, and past performance.
    \item Allocate available study hours across topics in priority order using a pure-function, testable scoring service.
    \item The LangGraph agent calls this service via tools rather than accessing the database directly, so it is structurally unable to bypass the deterministic business rules -- it can only explain the engine's output, not override it.
\end{itemize}

\subsection{Web architecture %
\BvrHint{\ldots web-based solution}}
A layered, one-directional web architecture:
\begin{itemize}[leftmargin=0.7cm]
    \item \textbf{Frontend (React):} responsive UI for subject/deadline entry, plan viewing, and progress logging.
    \item \textbf{Backend API (FastAPI):} routes $\rightarrow$ services $\rightarrow$ repositories $\rightarrow$ database, a strict one-directional layering enforced structurally rather than by convention.
    \item \textbf{Database (PostgreSQL, via SQLAlchemy ORM):} stores students, subjects, topics, plans, progress logs, and performance history.
\end{itemize}

\subsection{CI/CD and fast delivery enabler}
CI/CD will be used to:
\begin{itemize}[leftmargin=0.7cm]
    \item Automatically build and run tests on every change.
    \item Produce repeatable deployment artifacts to staging.
    \item Enable safe and frequent releases of incremental features.
\end{itemize}

\section{Evaluation Criterion %
\BvrHint{\ldots measurable and attributable}}
We define success using measurable metrics tied to the core adaptive
planning workflow.

\subsection{Primary evaluation metric}
\textbf{Plan adherence:} the percentage of planned study hours/sessions actually completed by the student, as logged in the system.
\begin{itemize}[leftmargin=0.7cm]
    \item Target: median plan adherence $\ge$ 70\% during pilot window.
    \item Attribution: measured from database timestamps comparing planned sessions vs. logged completions.
\end{itemize}

\subsection{Secondary evaluation metrics}
\begin{itemize}[leftmargin=0.7cm]
    \item \textbf{Rebalancing responsiveness:} how accurately and quickly the engine redistributes remaining workload after a missed or completed session, measured against remaining time-to-deadline.
    \item \textbf{Deadline miss rate:} percentage of tracked deadlines missed compared to a baseline (self-reported prior habits).
    \item \textbf{User trust/clarity:} student-reported clarity of plan reasoning using a simple 1--5 feedback question after viewing an agent explanation.
    \item \textbf{Reliability:} availability of staging/production for login and plan generation during typical study hours (e.g., $\ge$ 99\% uptime in pilot).
\end{itemize}

\subsection{Pilot validation plan %
\BvrHint{\ldots high-level}}
\begin{itemize}[leftmargin=0.7cm]
    \item Recruit 10--20 pilot student users across varying course loads.
    \item Compare plan adherence and self-reported stress/deadline-miss rate before/after within the iteration window.
\end{itemize}

\section{Scalability %
\BvrHint{\ldots with foresight}}
The proposition should scale in both theoretical and practical
senses.

\begin{enumerate}[leftmargin=0.7cm]
    \item \textbf{Scalable (theoretically):} The system should support growth in number of students, subjects, and concurrent plan re-computation by:
    \begin{itemize}[leftmargin=0.7cm]
        \item decoupling components (frontend/backend/agent),
        \item enabling pagination and indexing for common queries,
        \item using stateless API design for the deterministic engine.
    \end{itemize}

    \item \textbf{Operationally deployable (practically):} The system should run on common web hosting:
    \begin{itemize}[leftmargin=0.7cm]
        \item managed PostgreSQL database,
        \item containerized or platform-hosted deployment,
        \item CI/CD pipeline for staging/production.
    \end{itemize}
\end{enumerate}

\section{Engine availability heuristic}
A mature set of ``engines'' (tooling/services) is available to
implement this as a web-based project:
\begin{itemize}[leftmargin=0.7cm]
    \item Web framework for REST APIs (FastAPI) and reactive frontend (React).
    \item Managed relational database (PostgreSQL) with a mature ORM (SQLAlchemy).
    \item Agent orchestration framework (LangGraph) for the explanation layer.
    \item Authentication approach (token-based).
    \item Test framework for unit/integration tests (pytest).
    \item CI runner and staging deployment target.
\end{itemize}
We will use these mature components to focus on workload-engine
correctness, layered architecture, and measurement rather than
inventing new infrastructure.

\section{Project scope and deliverables %
\BvrHint{\ldots iteration-friendly}}
\subsection{Initial deliverable %
\BvrHint{\ldots first iteration}}
\begin{itemize}[leftmargin=0.7cm]
    \item Subject/topic/deadline input form + student profile
    \item Deterministic priority-scoring workload engine (pure functions)
    \item Generated study plan view
    \item Basic progress logging and timestamps for adherence measurement
    \item CI: build + backend unit tests + linting
    \item CD to staging with a single command or pipeline job
\end{itemize}

\subsection{Subsequent deliverables}
\begin{itemize}[leftmargin=0.7cm]
    \item Continuous feedback loop: re-balancing plan from logged progress
    \item LangGraph agent for plan explanation via tool calls
    \item Performance-history-aware prioritization
    \item Progress dashboard with adherence and rebalancing metrics
\end{itemize}

\section{Risks and mitigations}
\begin{itemize}[leftmargin=0.7cm]
    \item \textbf{Data privacy / consent:} minimize collected personal data; store only what is needed for planning and evaluation.
    \item \textbf{Student adoption:} keep input workflow minimal; provide sensible defaults so a first plan can be generated quickly.
    \item \textbf{Evaluation measurement ambiguity:} instrument timestamps and define clear session/completion states before development begins.
    \item \textbf{Agent overreach:} enforce strict tool-based access so the agent cannot bypass the deterministic engine's business rules.
\end{itemize}

\section{Summary}
This project proposes a deployable web platform that generates and
continuously adapts a personalized student study plan. It meets the
course criteria by emphasizing time-to-value through rapid
iterations, implementing CI/CD to support frequent safe releases,
and using measurable evaluation criteria (especially plan
adherence). It remains focused on engineering workflow, layered
architecture, and structural safety of the AI reasoning layer rather
than research-driven novelty.

\end{document}